% !PRINTED 510
\weberpage{510}

3. \gesperrt{If $\alpha$ is divisible by $\beta$, and $\beta$ is divisible by
$\gamma$, then $\alpha$ is also divisible by $\gamma$.}

For according to the assumption there are two integral functionals $\varkappa$,
$\lambda$ which satisfy the conditions $\alpha = \varkappa\beta$, $\beta = \lambda\gamma$.
Accordingly we have also $\alpha = \varkappa\lambda\gamma$, and since $\varkappa\lambda$
is integral, $\alpha$ is divisible by $\gamma$.

A product $\beta\gamma$ of two integral functionals is divisible both by $\beta$ and by
$\gamma$, and consequently, if $\beta$ is divisible by $\alpha$, then $\beta\gamma$ is
also divisible by $\alpha$. From these theorems it follows:

4. \gesperrt{If $\alpha$, $\beta$, \ldots\ are divisible by $\delta$, and if $\xi$,
$\eta$, \ldots\ are arbitrary integral functionals, then $\xi\alpha + \eta\beta + \ldots$
is divisible by $\delta$.}

These definitions and theorems extend, of course, also to the case in which numbers
take the place of functionals, and thus we obtain the divisibility of
integers.

5. \gesperrt{Two integral functionals $\alpha$, $\beta$ which are mutually divisible
by one another are called associated.}

6. \gesperrt{An integral functional $\varepsilon$ associated with the natural number 1,
i.\,e. every divisor of the number 1, is called a unit.}

According as $\varepsilon$ is a functional or a number, $\varepsilon$ is a
\gesperrt{functional} or a \gesperrt{numerical} unit.

In the field $R$ of the rational numbers, the primitive integral functions and the
quotients of two among them are to be regarded as the functional units. Of numerical
units there are in $R$ only two, namely $+\,1$ and $-\,1$.

Concerning the concepts hereby introduced, which stand in close mutual relation,
we derive a series of theorems.

7. \gesperrt{If $\alpha$, $\beta$ are associated functionals, then the quotients
$\beta : \alpha$ and $\alpha : \beta$ are units.}

For if we set $\beta = \alpha\varepsilon$, then $\varepsilon$ is an integral
functional and $1 : \varepsilon = \alpha : \beta$ is likewise integral; hence
$\varepsilon$ is a divisor of the number 1, i.\,e. $\varepsilon$ is a unit.

8. \gesperrt{If $\alpha$ is an integral functional and $\varepsilon$ a unit, then
$\alpha$ and $\alpha\varepsilon$ are associated.}

This follows immediately from the definition; for $\alpha : \alpha\varepsilon =
1 : \varepsilon$ and $\alpha\varepsilon : \alpha = \varepsilon$ are both integral
functionals.
